\documentclass[10pt,a4paper]{amsart}
\usepackage{amsmath,amssymb,amsthm,graphicx}
\usepackage{verbatim}
\newenvironment{reference}{\comment}{\endcomment}
%\newenvironment{reference}{}{}
\newenvironment{slogan}{\comment}{\endcomment}
\newenvironment{history}{\comment}{\endcomment}

% For commutative diagrams we use Xy-pic
\usepackage[all]{xy}

% We use 2cell for 2-commutative diagrams.
\xyoption{2cell}
% Theorem environments.
%
\theoremstyle{plain}
\newtheorem{theorem}[subsection]{Theorem}
\newtheorem{proposition}[subsection]{Proposition}
\newtheorem{lemma}[subsection]{Lemma}

\theoremstyle{definition}
\newtheorem{definition}[subsection]{Definition}
\newtheorem{example}[subsection]{Example}
\newtheorem{exercise}[subsection]{Exercise}
\newtheorem{situation}[subsection]{Situation}

\theoremstyle{remark}
\newtheorem{remark}[subsection]{Remark}
\newtheorem{remarks}[subsection]{Remarks}

\numberwithin{equation}{subsection}

% Macros
%
\def\lim{\mathop{\mathrm{lim}}\nolimits}
\def\colim{\mathop{\mathrm{colim}}\nolimits}
\def\Spec{\mathop{\mathrm{Spec}}}
\def\Hom{\mathop{\mathrm{Hom}}\nolimits}
\def\Ext{\mathop{\mathrm{Ext}}\nolimits}
\def\SheafHom{\mathop{\mathcal{H}\!\mathit{om}}\nolimits}
\def\SheafExt{\mathop{\mathcal{E}\!\mathit{xt}}\nolimits}
\def\Sch{\mathit{Sch}}
\def\Mor{\mathop{\mathrm{Mor}}\nolimits}
\def\Ob{\mathop{\mathrm{Ob}}\nolimits}
\def\Sh{\mathop{\mathit{Sh}}\nolimits}
\def\NL{\mathop{N\!L}\nolimits}
\def\CH{\mathop{\mathrm{CH}}\nolimits}
\def\proetale{{pro\text{-}\acute{e}tale}}
\def\etale{{\acute{e}tale}}
\def\QCoh{\mathit{QCoh}}
\def\Ker{\mathop{\mathrm{Ker}}}
\def\Im{\mathop{\mathrm{Im}}}
\def\Coker{\mathop{\mathrm{Coker}}}
\def\Coim{\mathop{\mathrm{Coim}}}

% Boxtimes
%
\DeclareMathSymbol{\boxtimes}{\mathbin}{AMSa}{"02}

%
% Macros for moduli stacks/spaces
%
\def\QCohstack{\mathcal{QC}\!\mathit{oh}}
\def\Cohstack{\mathcal{C}\!\mathit{oh}}
\def\Spacesstack{\mathcal{S}\!\mathit{paces}}
\def\Quotfunctor{\mathrm{Quot}}
\def\Hilbfunctor{\mathrm{Hilb}}
\def\Curvesstack{\mathcal{C}\!\mathit{urves}}
\def\Polarizedstack{\mathcal{P}\!\mathit{olarized}}
\def\Complexesstack{\mathcal{C}\!\mathit{omplexes}}
% \Pic is the operator that assigns to X its picard group, usage \Pic(X)
% \Picardstack_{X/B} denotes the Picard stack of X over B
% \Picardfunctor_{X/B} denotes the Picard functor of X over B
\def\Pic{\mathop{\mathrm{Pic}}\nolimits}
\def\Picardstack{\mathcal{P}\!\mathit{ic}}
\def\Picardfunctor{\mathrm{Pic}}
\def\Deformationcategory{\mathcal{D}\!\mathit{ef}}
\usepackage{cleveref}
\crefname{originallemma}{original lemma}{original lemmas}
\Crefname{originallemma}{Original lemma}{Original lemmas}
\makeatletter
\newlabel{algebra-lemma-dominate}{{00IA}{source}{}{sourcepointer-1}{}}
\newlabel{algebra-lemma-dominate@cref}{{[originallemma][0][]00IA}{[1][0][]source}{}{}{}}
\makeatother
\crefname{originallemma}{original lemma}{original lemmas}
\Crefname{originallemma}{Original lemma}{Original lemmas}
\makeatletter
\newlabel{algebra-lemma-ff-rings}{{00HQ}{source}{}{sourcepointer-2}{}}
\newlabel{algebra-lemma-ff-rings@cref}{{[originallemma][0][]00HQ}{[1][0][]source}{}{}{}}
\makeatother
\crefname{originallemma}{original lemma}{original lemmas}
\Crefname{originallemma}{Original lemma}{Original lemmas}
\makeatletter
\newlabel{lemma-base-change-flat}{{03MO}{source}{}{sourcepointer-3}{}}
\newlabel{lemma-base-change-flat@cref}{{[originallemma][0][]03MO}{[1][0][]source}{}{}{}}
\makeatother
\crefname{originaldefinition}{original definition}{original definitions}
\Crefname{originaldefinition}{Original definition}{Original definitions}
\makeatletter
\newlabel{more-algebra-definition-extension-valuation-rings}{{0ASG}{source}{}{sourcepointer-4}{}}
\newlabel{more-algebra-definition-extension-valuation-rings@cref}{{[originaldefinition][0][]0ASG}{[1][0][]source}{}{}{}}
\makeatother
\crefname{originallemma}{original lemma}{original lemmas}
\Crefname{originallemma}{Original lemma}{Original lemmas}
\makeatletter
\newlabel{spaces-properties-lemma-points-cartesian}{{03H4}{source}{}{sourcepointer-5}{}}
\newlabel{spaces-properties-lemma-points-cartesian@cref}{{[originallemma][0][]03H4}{[1][0][]source}{}{}{}}
\makeatother
\title{Lifting a valuation ring through a flat morphism of algebraic spaces}
\author{The Stacks Project Authors}
\date{stacks-project:spaces-morphisms.tex:local-snapshot-20261012; GNU FDL 1.2 or later}
\begin{document}
\maketitle
\setcounter{subsection}{1}
\noindent\textit{Editorial note.} All delivered mathematical text is the author's own, copied from the pinned Stacks Project chapter file spaces-morphisms.tex; no statement or proof step was rewritten, reordered or removed.\par\medskip
\noindent\textit{Editorial note.} The proof is the flat going-down argument written out in full; the hypothesis that the closed point lies in the image of |X| -> |Y| is used exactly once, to place the point u.\par\medskip
\noindent\textit{Editorial note.} The five references of the counted spans are recorded as non-load-bearing navigation pointers; the argument that consumes them stays inside the counted proof span.\par\medskip
\noindent\textit{Editorial note.} Preamble: the three native excerpt spans of sources/stacks-spaces-morphisms/preamble.tex (lines 8-13, 14-18 and 154-217). The chapter file itself carries no preamble text.\par\medskip
\noindent\textit{Allowed standard prerequisites (editorial).}\begin{itemize}
\item Description of points of a fibre product of algebraic spaces. Properties of Spaces, Lemma spaces-properties-lemma-points-cartesian, used at proof lines 8977-8980. Lifts the closed point of Spec(A) to a point of X\_A.
\item Flatness is stable under base change. Lemma lemma-base-change-flat, used at proof lines 8977-8978. Keeps flatness of X\_A -> Spec(A).
\item The going-down property of flat local ring maps: primes lying over the generic point. Algebra, Lemma algebra-lemma-ff-rings, used at proof lines 8983-8987. Produces the prime q of O\_\{U,u\} lying over (0) in A.
\item Existence of a valuation ring dominating a local domain inside its fraction field. Algebra, Lemma algebra-lemma-dominate, used at proof lines 8988-8990. Produces A' contained in kappa(q) dominating B/q.
\item Extension of valuation rings. More on Algebra, Definition more-algebra-definition-extension-valuation-rings, used at proof lines 8971-8973. Names the resulting relation between A and A' in the statement.
\end{itemize}\medskip
\ifdefined\hypertarget\hypertarget{sourcepointer-1}{}\fi
\noindent\textit{Source reference.} Original lemma 00IA: focused-seeds/AGP/sources/stacks-snapshot-20261011/tags.txt, source lines 672--672. Navigation pointer into the Stacks Project; the standard fact it supplies is applied inside the counted proof span and the argument resting on it stays there.\par\medskip
\ifdefined\hypertarget\hypertarget{sourcepointer-2}{}\fi
\noindent\textit{Source reference.} Original lemma 00HQ: focused-seeds/AGP/sources/stacks-snapshot-20261011/tags.txt, source lines 652--652. Navigation pointer into the Stacks Project; the standard fact it supplies is applied inside the counted proof span and the argument resting on it stays there.\par\medskip
\ifdefined\hypertarget\hypertarget{sourcepointer-3}{}\fi
\noindent\textit{Source reference.} Original lemma 03MO: focused-seeds/AGP/sources/stacks-snapshot-20261011/tags.txt, source lines 4756--4756. Navigation pointer into the Stacks Project; the standard fact it supplies is applied inside the counted proof span and the argument resting on it stays there.\par\medskip
\ifdefined\hypertarget\hypertarget{sourcepointer-4}{}\fi
\noindent\textit{Source reference.} Original definition 0ASG: focused-seeds/AGP/sources/stacks-snapshot-20261011/tags.txt, source lines 13579--13579. Navigation pointer into the Stacks Project; the standard fact it supplies is applied inside the counted proof span and the argument resting on it stays there.\par\medskip
\ifdefined\hypertarget\hypertarget{sourcepointer-5}{}\fi
\noindent\textit{Source reference.} Original lemma 03H4: focused-seeds/AGP/sources/stacks-snapshot-20261011/tags.txt, source lines 4553--4553. Navigation pointer into the Stacks Project; the standard fact it supplies is applied inside the counted proof span and the argument resting on it stays there.\par\medskip
\begin{lemma}
\label{lemma-lift-valuation-ring-through-flat-morphism}
Let $S$ be a scheme. Let $f : X \to Y$ be a flat morphism
of algebraic spaces over $S$.
Let $\Spec(A) \to Y$ be a morphism where $A$ is a
valuation ring. If the closed point of $\Spec(A)$ maps to a
point of $|Y|$ in the image of $|X| \to |Y|$, then there exists
a commutative diagram
$$
\xymatrix{
\Spec(A') \ar[r] \ar[d] & X \ar[d] \\
\Spec(A) \ar[r] & Y
}
$$
where $A \to A'$ is an extension of valuation rings
(More on Algebra, Definition
\ref{more-algebra-definition-extension-valuation-rings}).
\end{lemma}

\begin{proof}
The base change $X_A \to \Spec(A)$ is flat
(Lemma \ref{lemma-base-change-flat}) and the closed point of
$\Spec(A)$ is in the image of $|X_A| \to |\Spec(A)|$
(Properties of Spaces, Lemma \ref{spaces-properties-lemma-points-cartesian}).
Thus we may assume $Y = \Spec(A)$. Let $U \to X$ be a surjective
\'etale morphism where $U$ is a scheme. Let $u \in U$ map to
the closed point of $\Spec(A)$. Consider the flat local ring map
$A \to B = \mathcal{O}_{U, u}$. By
Algebra, Lemma \ref{algebra-lemma-ff-rings}
there exists a prime ideal $\mathfrak q \subset B$ such that
$\mathfrak q$ lies over $(0) \subset A$. By
Algebra, Lemma \ref{algebra-lemma-dominate}
we can find a valuation ring $A' \subset \kappa(\mathfrak q)$
dominating $B/\mathfrak q$. The induced morphism
$\Spec(A') \to U \to X$ is a solution to the problem
posed by the lemma.
\end{proof}

\end{document}
